\subsection{实验结论}
实验结果证明,即使是简单调整矩阵乘法内层访问顺序,使得对所有矩阵都是内存连续式访问,就能极大增加程序运行效率,充分体现了空间局部性对程序性能的深刻影响。
各种优化方法中循环重排、多线程并行是提升矩阵乘法性能的核心手段,其中取得最佳优化结果的就是6线程并行优化。在多线程应用中，应根据计算规模选择合适的线程数，
避免因过载导致的“过饱和”性能损耗,使得多线程的优势被线程切换等的损耗给掩盖。同时，QEMU 仿真环境能有效验证算法逻辑和相对趋势，但其绝对性能指标受限于软件翻译和宿主机核心资源。